% GENERATED FILE. Edit canonical lessons, then run npm run generate:books.
% canonical-source-hash: fnv1a64:c1c41bf5052885f9
% canonical-lessons: MR-C404-valsa, MR-C404-rahdari, MR-C404-signal, MR-C404-khetne, MR-C404-datne

\chapter{Traffic and Detours}
\label{ch:mr-traffic-and-detours}
\hlchaptermodality{404}

\begin{chapteropening}
\textbf{By the end of this chapter:} \emph{I can talk about traffic, signals, a detour and a crowd pressing in.}

Complete the last lesson of chapter 404: I can talk about traffic, signals, a detour and a crowd pressing in.
\end{chapteropening}

\section[\texorpdfstring{\mr{वळसा} (vaḷsā)}{वळसा (vaḷsā)}]{\mr{वळसा} (vaḷsā) — a detour; a roundabout way}

\label{lesson:MR-C404-valsa}



\begin{quote}
Before the new one: say the Marathi for to tip over, then the Marathi for to bang against.
\end{quote}

\subsection*{You'll want to know: \mr{वळसा}}
\textbf{\mr{वळसा}} — \emph{vaḷsā} — \textquotedblleft{}a detour; a roundabout way\textquotedblright{}.

\begin{quote}
\mr{पूल} \mr{बंद} \mr{होता} \mr{म्हणून} \mr{आम्ही} \mr{लांबचा} \mr{वळसा} \mr{घेतला}.
\end{quote}

\begin{quote}
\emph{pūl band hotā mhaṇūn āmhī lāmbcā vaḷsā ghetlā.}
\end{quote}

\begin{quote}
The bridge was closed, so we took a long detour.
\end{quote}

Masculine.

\mr{वळसा} \mr{घेणे} or \mr{वळसा} \mr{घालणे}, to go the long way round.

\begin{cousinweb}
\textbf{\mr{वळसा}} (\emph{vaḷsā}) is built on the stem of \textbf{\mr{वळणे}}, \textquotedblleft{}to turn\textquotedblright{}; how the noun was formed is not settled.
\end{cousinweb}

\begin{grammarlens}[title={What you've built}]
One more word for getting around.
\end{grammarlens}

\subsection*{Your turn}
\begin{itemize}
\raggedright
  \item \emph{Say it:} \emph{vaḷsā}
  \item \emph{Say it:} \emph{vaḷsā}, once more
  \item \emph{Say it:} \emph{pūl band hotā mhaṇūn āmhī lāmbcā vaḷsā ghetlā.} — The bridge was closed, so we took a long detour.
  \item \emph{Recall:} say the Marathi for to tip over, then the Marathi for to bang against, then say \emph{vaḷsā} again
\end{itemize}

\begin{tcolorbox}[breakable,colback=teal!4,colframe=teal!35!black,title={Before you move on}]
What does \mr{वळसा} mean? (\textquotedblleft{}A detour; a roundabout way\textquotedblright{}.) Say it once more.
\end{tcolorbox}
\section[\texorpdfstring{\mr{रहदारी} (rahdārī)}{रहदारी (rahdārī)}]{\mr{रहदारी} (rahdārī) — road traffic; movement}

\label{lesson:MR-C404-rahdari}



\begin{quote}
Before the new one: say the Marathi for to bang against, then the Marathi for a detour.
\end{quote}

\subsection*{You'll want to know: \mr{रहदारी}}
\textbf{\mr{रहदारी}} — \emph{rahdārī} — \textquotedblleft{}road traffic; movement\textquotedblright{}.

\begin{quote}
\mr{सकाळी} \mr{या} \mr{चौकात} \mr{खूप} \mr{रहदारी} \mr{असते}.
\end{quote}

\begin{quote}
\emph{sakāḷī yā caukāt khūp rahdārī aste.}
\end{quote}

\begin{quote}
In the morning there is a lot of traffic at this square.
\end{quote}

Feminine.

A feminine noun; \mr{वाहतूक} is the same idea in a more official word.

\begin{cousinweb}
\textbf{\mr{रहदारी}} (\emph{rahdārī}) is a loan usually traced to Persian \emph{rāhdārī}, built on \emph{rāh}, \textquotedblleft{}road\textquotedblright{}: the coming and going on a road.
\end{cousinweb}

\begin{grammarlens}[title={What you've built}]
One more word for getting around.
\end{grammarlens}

\subsection*{Your turn}
\begin{itemize}
\raggedright
  \item \emph{Say it:} \emph{rahdārī}
  \item \emph{Say it:} \emph{rahdārī}, once more
  \item \emph{Say it:} \emph{sakāḷī yā caukāt khūp rahdārī aste.} — In the morning there is a lot of traffic at this square.
  \item \emph{Recall:} say the Marathi for to bang against, then the Marathi for a detour, then say \emph{rahdārī} again
\end{itemize}

\begin{tcolorbox}[breakable,colback=teal!4,colframe=teal!35!black,title={Before you move on}]
What does \mr{रहदारी} mean? (\textquotedblleft{}Road traffic; movement\textquotedblright{}.) Say it once more.
\end{tcolorbox}
\section[\texorpdfstring{\mr{सिग्नल} (signal)}{सिग्नल (signal)}]{\mr{सिग्नल} (signal) — a traffic signal}

\label{lesson:MR-C404-signal}



\begin{quote}
Before the new one: say the Marathi for a detour, then the Marathi for road traffic.
\end{quote}

\subsection*{You'll want to know: \mr{सिग्नल}}
\textbf{\mr{सिग्नल}} — \emph{signal} — \textquotedblleft{}a traffic signal\textquotedblright{}.

\begin{quote}
\mr{सिग्नल} \mr{लाल} \mr{असला} \mr{की} \mr{गाड्या} \mr{थांबतात}.
\end{quote}

\begin{quote}
\emph{signal lāl aslā kī gāḍyā thāmbtāt.}
\end{quote}

\begin{quote}
When the signal is red, the cars stop.
\end{quote}

Masculine.

\mr{सिग्नल} \mr{तोडणे}, to jump a signal.

\begin{cousinweb}
\textbf{\mr{सिग्नल}} (\emph{signal}) is a loan from English \emph{signal}.
\end{cousinweb}

\begin{grammarlens}[title={What you've built}]
One more word for getting around.
\end{grammarlens}

\subsection*{Your turn}
\begin{itemize}
\raggedright
  \item \emph{Say it:} \emph{signal}
  \item \emph{Say it:} \emph{signal}, once more
  \item \emph{Say it:} \emph{signal lāl aslā kī gāḍyā thāmbtāt.} — When the signal is red, the cars stop.
  \item \emph{Recall:} say the Marathi for a detour, then the Marathi for road traffic, then say \emph{signal} again
\end{itemize}

\begin{tcolorbox}[breakable,colback=teal!4,colframe=teal!35!black,title={Before you move on}]
What does \mr{सिग्नल} mean? (\textquotedblleft{}A traffic signal\textquotedblright{}.) Say it once more.
\end{tcolorbox}
\section[\texorpdfstring{\mr{खेटणे} (kheṭṇe)}{खेटणे (kheṭṇe)}]{\mr{खेटणे} (kheṭṇe) — to brush against}

\label{lesson:MR-C404-khetne}



\begin{quote}
Before the new one: say the Marathi for road traffic, then the Marathi for a traffic signal.
\end{quote}

\subsection*{You'll want to know: \mr{खेटणे}}
\textbf{\mr{खेटणे}} — \emph{kheṭṇe} — \textquotedblleft{}to brush against\textquotedblright{}.

\begin{quote}
\mr{गर्दीत} \mr{एक} \mr{माणूस} \mr{मला} \mr{खेटून} \mr{गेला}.
\end{quote}

\begin{quote}
\emph{gardīt ek māṇūs malā kheṭūn gelā.}
\end{quote}

\begin{quote}
In the crowd a man brushed past me.
\end{quote}

Past (he …): \mr{तो} \mr{खेटला}, \emph{to kheṭlā}.

Often in the -\mr{ऊन} form: \mr{खेटून} \mr{उभं} \mr{राहणे}, to stand pressed close.

\begin{cousinweb}
\textbf{\mr{खेटणे}} (\emph{kheṭṇe}) is a Marathi verb for touching someone or something as you pass very close; its origin is not settled.
\end{cousinweb}

\begin{grammarlens}[title={What you've built}]
One more word for getting around.
\end{grammarlens}

\subsection*{Your turn}
\begin{itemize}
\raggedright
  \item \emph{Say it:} \emph{kheṭṇe}
  \item \emph{Say it:} \emph{kheṭṇe}, once more
  \item \emph{Say it:} \emph{gardīt ek māṇūs malā kheṭūn gelā.} — In the crowd a man brushed past me.
  \item \emph{Recall:} say the Marathi for road traffic, then the Marathi for a traffic signal, then say \emph{kheṭṇe} again
\end{itemize}

\begin{tcolorbox}[breakable,colback=teal!4,colframe=teal!35!black,title={Before you move on}]
What does \mr{खेटणे} mean? (\textquotedblleft{}To brush against\textquotedblright{}.) Say it once more.
\end{tcolorbox}
\section[\texorpdfstring{\mr{दाटणे} (dāṭṇe)}{दाटणे (dāṭṇe)}]{\mr{दाटणे} (dāṭṇe) — to crowd thickly}

\label{lesson:MR-C404-datne}



\begin{quote}
Before the new one: say the Marathi for a traffic signal, then the Marathi for to brush against.
\end{quote}

\subsection*{You'll want to know: \mr{दाटणे}}
\textbf{\mr{दाटणे}} — \emph{dāṭṇe} — \textquotedblleft{}to crowd thickly\textquotedblright{}.

\begin{quote}
\mr{संध्याकाळी} \mr{रस्त्यावर} \mr{गर्दी} \mr{दाटली}.
\end{quote}

\begin{quote}
\emph{sandhyākāḷī rastyāvar gardī dāṭlī.}
\end{quote}

\begin{quote}
In the evening the crowd thickened on the road.
\end{quote}

Past (he …): \mr{तो} \mr{दाटला}, \emph{to dāṭlā}.

Clouds, crowds and feelings all \mr{दाटतात}: \mr{ढग} \mr{दाटले}, the clouds gathered.

\begin{cousinweb}
\textbf{\mr{दाटणे}} (\emph{dāṭṇe}) is inherited: CDIAL lists Marathi \emph{dāṭṇẽ}, \textquotedblleft{}to be pressed in, become thick\textquotedblright{}, under \emph{*draṭ} (CDIAL 6618).
\end{cousinweb}

\begin{grammarlens}[title={What you've built}]
One more word for getting around.
\end{grammarlens}

\subsection*{Your turn}
\begin{itemize}
\raggedright
  \item \emph{Say it:} \emph{dāṭṇe}
  \item \emph{Say it:} \emph{dāṭṇe}, once more
  \item \emph{Say it:} \emph{sandhyākāḷī rastyāvar gardī dāṭlī.} — In the evening the crowd thickened on the road.
  \item \emph{Recall:} say the Marathi for a traffic signal, then the Marathi for to brush against, then say \emph{dāṭṇe} again
\end{itemize}

\begin{tcolorbox}[breakable,colback=teal!4,colframe=teal!35!black,title={Before you move on}]
What does \mr{दाटणे} mean? (\textquotedblleft{}To crowd thickly\textquotedblright{}.) Say it once more.
\end{tcolorbox}
